\hypertarget{sec:planteamiento}{}%
\begin{center}
    {\Large\bfseries 1. PLANTEAMIENTO DEL PROBLEMA}
\end{center}
\vspace{1em}

\hypertarget{sec:antecedentes}{}%
\noindent{\large\bfseries 1.1 ANTECEDENTES}\par\vspace{0.5em}

En los últimos años, el sector financiero global y nacional ha experimentado una acelerada transformación digital impulsada por el desarrollo de la tecnología financiera (fintech) y la masificación de los dispositivos móviles \cite{bermeo2022transformacion}. En Colombia, plataformas de servicios financieros digitales y billeteras electrónicas como Nequi, DaviPlata y las aplicaciones móviles de la banca tradicional (destacando Bancolombia) se han consolidado como herramientas fundamentales para la inclusión financiera, transformando radicalmente la forma en que las personas gestionan sus recursos monetarios, realizan transferencias y ejecutan pagos cotidianos sin necesidad de recurrir al efectivo \cite{banca2024reporte}.

Estudios sobre adopción tecnológica en jóvenes universitarios demuestran que este grupo poblacional es uno de los más receptivos y dependientes de las soluciones de banca móvil \cite{lopez2023adopcion}. Las dinámicas académicas, caracterizadas por transacciones inmediatas de bajo monto —como el pago de fotocopias, alimentación en cafeterías, transporte urbano y transferencias informales entre pares—, encuentran en las aplicaciones financieras un medio ágil que optimiza el tiempo y elimina las barreras asociadas al manejo de dinero físico o la asistencia a sucursales bancarias \cite{restrepo2023habitos}.

Sin embargo, el funcionamiento de estas plataformas digitales en entornos universitarios no ocurre en un plano aislado, sino que está estrechamente condicionado por la estabilidad de su infraestructura tecnológica, la disponibilidad de los servidores y la calidad de la conectividad \cite{martinez2023retos}. A diferencia de los centros urbanos con alta densidad comercial e infraestructura bancaria diversificada, las comunidades universitarias integradas en ciudades intermedias como Tunja presentan patrones de uso altamente concentrados. En el campus de la Universidad Pedagógica y Tecnológica de Colombia (UPTC), Sede Tunja, la masa estudiantil recurre de manera simultánea a las mismas plataformas en franjas horarias específicas, lo que genera una alta vulnerabilidad ante eventuales caídas del servicio o fallas operativas en los sistemas de las entidades \cite{pulido2024caracteristicas}.

Investigaciones previas sobre la adopción de tecnologías financieras suelen enfocar sus análisis en indicadores macroeconómicos de bancarización o en la usabilidad técnica de las interfaces de usuario desde la perspectiva del desarrollador de software \cite{silva2022evaluacion}. No obstante, existe una escasa documentación científica y empírica a nivel local que evalúe cuantitativa y cualitativamente el grado de dependencia funcional que han desarrollado los estudiantes universitarios frente a aplicaciones específicas como Nequi, DaviPlata y Bancolombia, así como la magnitud del impacto o parálisis en las actividades cotidianas que ocasiona la interrupción de estos servicios dentro del campus.

Esta brecha de conocimiento entre la penetración tecnológica de las billeteras digitales y la capacidad de contingencia real de los estudiantes universitarios ante la indisponibilidad del servicio constituye el punto de partida del presente estudio en la UPTC, Sede Tunja.

\vspace{1em}

\hypertarget{sec:descripcion}{}%
\noindent{\large\bfseries 1.2 DESCRIPCIÓN DE LA SITUACIÓN PROBLEMÁTICA}\par
En la cotidianidad de la Universidad Pedagógica Y Tecnológica de Colombia (UPTC), Sede Central Tunja, el uso de billeteras digitales e intermediarios financieros móviles como Nequi, DaviPlata y la aplicación de Bancolombia se ha integrado de forma generalizada en las dinámicas de consumo y movilidad de la comunidad estudiantil. La baja tenencia de dinero en efectivo entre los jóvenes universitarios ha provocado que transacciones esenciales tales como la compra de alimentos en las cafeterías del campus, el pago de fotocopias e impresiones de trabajos académicos, el pago del transporte público o colectivo urbano, e incluso la cancelación de arriendos en las zonas residenciales aledañas dependan casi de manera absoluta de la disponibilidad y correcto funcionamiento de estos sistemas informáticos.

A pesar de la alta conveniencia y agilidad que estas plataformas ofrecen, su infraestructura tecnológica presenta problemas operativos manifestados en caídas frecuentes de plataforma, interrupciones no programadas en los servidores, fallas de procesamiento durante picos de tráfico y congelamiento de saldos o transacciones pendientes. Cuando alguna de estas aplicaciones deja de operar o experimenta latencias prolongadas, el estudiante queda expuesto a un escenario de parálisis financiera inmediata. En el contexto de la UPTC Sede Tunja, esta situación genera problemas concretos como la imposibilidad de acceder a servicios básicos de alimentación durante los recesos de clase, retrasos en la entrega de actividades académicas físicas por no poder pagar la impresión de documentos, o dificultades de desplazamiento al no contar con efectivo para abordar el transporte urbano.


\vspace{1em}

\hypertarget{sec:formulacion}{}%
\noindent{\large\bfseries 1.3 FORMULACIÓN DEL PROBLEMA}\par
A partir de la situación problemática descrita, la pregunta general de investigación que orienta el presente estudio es:

¿De qué manera la disponibilidad y el estado operativo de las aplicaciones financieras digitales Nequi, DaviPlata y Bancolombia condicionan  las actividades cotidianas (académicas, de alimentación y movilidad) e inciden en el grado de dependencia de los estudiantes universitarios de la UPTC, Sede Tunja?

Para abordar de manera sistemática la pregunta central, la investigación se desglosa en las siguientes sub-preguntas:

1.	¿Cuál es el nivel de dependencia de los estudiantes frente a cada una de estas aplicaciones y cuáles son las actividades cotidianas más afectadas cuando se presentan interrupciones en el servicio?

2.	¿Existe un contraste significativo en el grado de vitalidad e impacto operacional entre Nequi, DaviPlata y Bancolombia durante una falla o caída de sistema?

3.	¿Qué mecanismos de contingencia (como uso de efectivo de reserva, migración a plataformas alternativas o cancelación de compras) implementan los estudiantes de la UPTC ante la indisponibilidad de su billetera digital preferida?


\vspace{1em}

\hypertarget{sec:consecuencias}{}%
\noindent{\large\bfseries 1.4 CONSECUENCIAS DE NO INVESTIGAR}\par
De no realizarse esta investigación, continuara la parálisis que sufren los estudiantes de la UPTC ante las caídas de Nequi, DaviPlata y Bancolombia, invisibilizando la alta vulnerabilidad a la que se ven expuestos en sus actividades diarias de alimentación, transporte y cumplimiento académico. Asimismo no permitirá el desarrollo de los comercios del campus y a las entidades bancarias de insumos fundamentados para promover alternativas de pago o exigir mayores estándares de disponibilidad en plataformas financieras digitales.

\vspace{1em}

\hypertarget{sec:delimitacion}{}%
\noindent{\large\bfseries 1.5 DELIMITACIÓN}\par

La presente investigación se enfoca en caracterizar el nivel de dependencia hacia las billeteras digitales (Nequi, DaviPlata y Bancolombia) y evaluar el impacto que generan sus caídas operativas en las actividades cotidianas de los estudiantes de la UPTC, Sede Tunja. En el Cuadro 1 se detallan las dimensiones de la delimitación aplicada al estudio.

\begin{table}[htbp]
\centering
\begin{tabular}{|p{3cm}|p{12cm}|}
\hline
\textbf{Dimensión} & \textbf{Delimitación aplicada} \\ \hline
Temática & Análisis de la dependencia funcional e impacto operacional causado por fallas en tres aplicaciones de banca móvil (Nequi, DaviPlata y Bancolombia) en actividades cotidianas (alimentación, transporte, pagos académicos); no abarca el desarrollo de software ni auditorías de seguridad informática. \\ \hline
Poblacional & Muestra de estudiantes universitarios de pregrado matriculados en la UPTC, Sede Tunja, que utilicen activamente al menos una de las tres plataformas analizadas. \\ \hline
Espacial & Campus central de la Universidad Pedagógica y Tecnológica de Colombia (UPTC), Sede Tunja, incluyendo comercios internos (cafeterías, centros de fotocopiado) y zonas residenciales aledañas. \\ \hline
Temporal & Marco temporal de 16 semanas, correspondiente al semestre lectivo del año 2026. \\ \hline
\end{tabular}
\caption{Delimitación de la investigación}
\label{tab:delimitacion}
\end{table}

\begin{table}[htbp]
\centering
\begin{tabular}{|p{4.5cm}|p{10.5cm}|}
\hline
\textbf{Queda fuera} & \textbf{Por qué} \\ \hline
Otras aplicaciones financieras o criptomonedas & Se evalúan únicamente Nequi, DaviPlata y Bancolombia por ser las de mayor penetración y uso dentro de la comunidad universitaria local. \\ \hline
Estudiantes de otras sedes (Sogamoso, Duitama, Chiquinquirá) & Las dinámicas de movilidad, acceso comercial y uso del campus varían según la sede; el estudio se concentra en Tunja. \\ \hline
Análisis del código de programación de las aplicaciones & El estudio aborda la perspectiva del usuario y el impacto operativo, no el desarrollo o mantenimiento del software bancario. \\ \hline
\end{tabular}
\caption{Criterios de exclusión del estudio}
\label{tab:exclusion}
\end{table}
